\documentclass[10pt,letterpaper]{article}
\usepackage{cornell-notes}

\title{Clause 28.07.06.19: Riemann Zeta Function}
\author{}
\date{}
\setDocTitle{Clause 28.07.06.19: Riemann Zeta Function}
\setDocSubtitle{C++ 2024}
\setDocOwner{C++ 2024 Cornell Notes}
\setDocDate{\today}
\setCornellCollection{C++ 2024 Cornell Notes}
\setCornellUnitType{Clause}
\setCornellUnitNumber{28.07.06.19}
\setCornellUnitTitle{Riemann Zeta Function}
\hypersetup{pdftitle={Clause 28.07.06.19: Riemann Zeta Function},pdfsubject={Cornell notes},pdfauthor={}}

\begin{document}
\maketitle
\begin{CornellOverviewBox}{Overview}
\textbf{Classification:} Standard library - Normative\\[2pt]
\textbf{Learning objective:} Understand the rules, terminology, and semantic consequences associated with Riemann zeta function.
\end{CornellOverviewBox}\begin{CornellSummaryBox}{Summary}
{\sffamily\bfseries\color{Accent} Summary}\par\vspace{3pt}Section 28.7.6.19 focuses on Riemann zeta function. Apply the specified interface together with its mandates, effects, result conditions, exception behavior, and complexity constraints. When applying it, preserve the distinction between mandatory requirements, recommendations, permissions, and behavior deliberately left undefined, unspecified, or implementation-defined.
\end{CornellSummaryBox}

\end{document}
